\section{Caps and two-point uniqueness}
\label{sec:caps}

The next two sections supply the local input for the harmonic-height
argument.  A cap controls the interaction with generators outside the
operated colors.  On a spherical coset, this control gives both a convex
description of a layer and a uniqueness theorem for its possible minima.
We retain the family $E_L(x)=B_xP_{f,L}$ throughout.  In particular, no
adjacency assumption between $f$ and the operated colors is made.
The two-point conclusion must identify the actual coset, rather than
only its numerical layer data.  We first build convex subobjects of a
capped layer and then use right simple modules to recover that coset.

\subsection{Accessibility of a cap}

\begin{definition}[Cap]
\label{cap:definition}
For $q\in\mathbb Z$ and a color set $D$, the family at $x$ is
\emph{capped at $q$ on $D$} if
\[
 m_i^d(x)=0\qquad(i\in D,\ d>q).
\]
Equivalently, every family member has no generator in those colors
above level $q$.  We write $\partial J$ for the exterior neighbors of
$J$ in $\widehat S$.
\end{definition}

Positivity of the family weights proves the equivalence in the
definition.  By Proposition~\ref{cat:family}, data outside $J$ are
fixed on $A_Jx$, so a cap on $\partial J$ is independent of its
representative.

\begin{lemma}[A capped negative operation]
\label{cap:capped-negative}
Let $E$ be a bounded projective complex capped at $q$ on
$J\cup\partial J$, and let $G$ be any word of negative twists in $J$.
Then $GE$ has the same cap, and
\begin{equation}
 Y^q(GE)\simeq Y^q\bigl(G(Y^q(E))\bigr).
 \label{cap:negative-layer}
\end{equation}
\end{lemma}
\begin{proof}
The high quotient $R=\tau_{\ge q+1}E$ is supported outside
$J\cup\partial J$.  Its tests at $J$ vanish termwise, so every
evaluation or coevaluation term in $J$ vanishes and $GR\simeq R$.
Apply $G$ to the truncation triangle
$\tau_{\le q}E\to E\to R$.  Negative twists do not increase levels.
The low term therefore remains at levels at most $q$, while the
unchanged quotient is at levels at least $q+1$.  The long exact
layer sequence identifies all higher layers with those of $R$ and
identifies level $q$ with that of $G\tau_{\le q}E$.  Apply $G$ next
to $\tau_{\le q-1}E\to\tau_{\le q}E\to Y^q(E)$.
Its first term stays below $q$, proving \eqref{cap:negative-layer}.
\end{proof}

\begin{lemma}[Cap accessibility]
\label{cap:access}
Let $J\subseteq S$ be spherical and suppose $A_Jx$ has a cap at $q$
on $\partial J$.  It has a representative capped at $q$ also on $J$.
Every negative letter in $J$ preserves this joint cap.  A positive
letter in $J$ preserves it whenever the top at that color in the
level-$q$ layer vanishes in every family member.
\end{lemma}
\begin{proof}
Negative preservation is Lemma~\ref{cap:capped-negative}.  For the
positive assertion use the same high quotient $R$.  It is unchanged.
On the low truncation, the only possible level-$q+1$ output is the
upper part of the old level-$q$ layer.  By Lemma~\ref{lay:positive}
this is precisely its top in the operated color; it vanishes by
hypothesis.

For accessibility, let $d>q$ be the greatest occupied level on $J$.
The joint cap at $d$ holds.  At that level the portion on $J$ is a
nonzero direct summand of a pure layer, since the boundary is absent.
It has a nonzero top somewhere on $J$.  Apply a negative letter at
such a color and continue while level $d$ remains.

The selected colors form a reduced Coxeter expression.  Indeed,
suppose the letters already selected form a reduced word $w$, with
the last operation on the left.  If the next color $j$ were a left
descent, Matsumoto moves could put $j$ at the last-action end of a
reduced expression for $w$.  The corresponding negative words give
the same action.  Every intermediate negative operation preserves
the joint $d$-cap.  Immediately before the last inverse twist, the
tests in color $j$ vanish above $d+2$; after it they vanish at $d+2$
as well, by the negative test shift.  Lemma~\ref{lay:extremal} would
then make the current level-$d$ top at $j$ zero, contrary to its
selection.  Thus each choice is a Coxeter ascent.  The longest length
in $W_J$ bounds the number of choices before level $d$ disappears.
No larger level is recreated.  Repeating at successive excess
levels reaches $q$ in finitely many stages.  The empty set $J$
requires no operations.
\end{proof}

\subsection{Candidate cosets and the uniqueness problem}

Fix a nonempty spherical set $T\subseteq S$, an element $x_0\in A$,
and an integer $q$.  Suppose the family $(E_L(x_0))_L$ is capped at
$q$ on $\partial T$.  Exterior invariance preserves this cap
throughout $A_Tx_0$.  In the comparison below, all candidates refer to this
family, this coset and this level.

\begin{definition}[Candidate]
\label{two:definition}
A candidate is $(J,A_Jz)$ with $J\subsetneq T$ and $z\in A_Tx_0$.
Put $V=T\setminus J$, let $p$ harmonically interpolate $m^q(z)$
on $J$, and put $a=-(Cp)_T$.  The requirements are:
\begin{enumerate}
\item the open colors $V$ are capped at $q$;
\item $a$ is strictly positive on $V$ and zero on $J$;
\item every representative $y\in A_Jz$ capped at $q$ on $T$
satisfies $m_J^q(y)\ge p_J$.
\end{enumerate}
\end{definition}

These are properties of the typed coset: its exterior counts, hence
$p$ and $a$, are fixed.  The representatives in the third condition
exist by Lemma~\ref{cap:access}, since
$\partial J\subseteq V\cup\partial T$.

\begin{proposition}[Two-point uniqueness]
\label{two:uniqueness}
Let $T\subseteq S$ be nonempty and spherical, let $x_0\in A$, and let
$q\in\mathbb Z$ cap $(E_L(x_0))_L$ on $\partial T$.
There is at most one typed coset $(J,A_Jz)$ satisfying
Definition~\ref{two:definition} in $A_Tx_0$ at level $q$.
\end{proposition}

The proof first identifies the harmonic vectors of two candidates;
the positivity condition then identifies their open colors and types.
It must also show that their representatives differ by an element of
the common parabolic subgroup.  Capped-layer subobjects will supply
the numerical comparison, while transported right simple modules will
detect this relation between representatives.

\subsection{Convex subobjects of a capped layer}

We first express the excess of a capped layer over its harmonic
interpolation through dimensions of subobjects.  The weak form below
will apply to candidates; the strict form will also be needed in
Section~\ref{sec:null}.

For a spherical $J$, harmonic interpolation of exterior data is
unique: $p$ agrees with those data outside $J$ and satisfies
$(Cp)_J=0$.  It is nonnegative for nonnegative exterior data and is
fixed by $W_J$, by Lemma~\ref{cat:cartan}.

\begin{proposition}[The capped-layer hull]
\label{cap:hull}
Suppose $x$ is capped at $q$ on $J\cup\partial J$, with $J$ spherical.
Let $p$ be the harmonic interpolation of $m^q(x)$ on $J$.  Assume
\begin{equation}
 m_j^q(y)\ge p_j\quad(j\in J)
 \quad\text{for every jointly capped }y\in A_Jx.
 \label{cap:universal}
\end{equation}
Put $e=m^q(x)-p$, which is supported on $J$. Then $e_J$ belongs to
the convex hull of finitely many
vectors $s_J$ for which $s$ is supported on $J$ and there are graded
subrepresentations
\begin{equation}
 \mathcal S_L\subseteq Y^q(E_L(x)),\qquad
 \supp\mathcal S_L\subseteq\widetilde J,\qquad
 \sum_Lw_L\dim\mathcal S_L=\phi s.
 \label{cap:subobjects}
\end{equation}
If every inequality in \eqref{cap:universal} is strict, $e_J$ is in
the ordinary interior of such a hull in $\RR^J$.  For $J=\varnothing$
this uses the convention that $\{0\}$ is open in $\RR^0$.
\end{proposition}
\begin{proof}
For each $w\in W_J$ choose a negative reduced lift $G_w$.  Put
$M_L=Y^q(E_L(x))$.  The reduced-word amplitude in
Corollary~\ref{lay:reduced} places $G_wM_L$ at levels $q-1,q$.
Apply $G_w^{-1}$ to its low/high truncation triangle, obtaining
\[
 A_L\longrightarrow M_L\longrightarrow C_L.
\]
The endpoint bounds are $[q-1,q]$ and $[q,q+1]$.  The long exact
layer sequence forces the possible level $q-1$ of $A_L$ and level
$q+1$ of $C_L$ to vanish.  Thus both are pure at $q$ and
\begin{equation}
 0\longrightarrow A_L\longrightarrow M_L\longrightarrow C_L
       \longrightarrow0
 \label{cap:short-exact}
\end{equation}
is exact in the shifted heart.  Its subobject is an actual graded
subrepresentation by Lemma~\ref{lay:structure}.  The low part of
$G_wM_L$ is supported on $\widetilde J$, since outside-color counts
are unchanged, and applying $G_w^{-1}$ keeps that support.

The inclusion in \eqref{cap:short-exact} is obtained by twisting,
canonical truncation, and inverse twisting.  These operations commute
with label tensoring, including their morphisms into $M_L$.  The
weighted fusion identity therefore gives \eqref{cap:subobjects} for
a vector $s_w$.  This does not require compatible choices of Gaussian
eliminations among the family members.

The quotient has folded dimension $m^q(x)-s_w$.  Applying $G_w$
recovers the high part at level $q$, whose signed generator count is
therefore $w(m^q(x)-s_w)$.  Lemma~\ref{cap:capped-negative} identifies
it with the actual level-$q$ layer at the operated representative.
That representative is jointly capped, so \eqref{cap:universal}
and $wp=p$ give
\begin{equation}
 \bigl(w(e-s_w)\bigr)_J\ge0.
 \label{cap:chamber-inequality}
\end{equation}
In the strict case all these coordinates are strictly positive.

On the root-coordinate space supported on $J$, finite Coxeter
chambers cover the dual space.  Thus every real functional has the
form $\lambda(v)=\sum_{j\in J}c_j(wv)_j$ for some $w$ and $c_j\ge0$.
Equation \eqref{cap:chamber-inequality} gives
\[
 \lambda(e)\ge\lambda(s_w)
       \ge\min_{s\in\mathcal V}\lambda(s),
 \qquad \mathcal V=\{s_w:w\in W_J\}\cup\{0\}.
\]
Separation proves membership in this finite convex hull.  If the
inequalities are strict and $\lambda\ne0$, some $c_j>0$, so the first
inequality is strict.  A boundary point, including a point of a hull
in a proper affine subspace, would admit a nonzero supporting
functional with equality.  This proves ordinary interior membership.
For $J=\varnothing$ all the vectors are zero.
\end{proof}

The convex description permits a linear estimate proved for those
subobjects to be applied to the excess vector, even when that vector is
not itself the dimension vector of a subobject.

\subsection{Right simples, duality, and parabolic detection}

Fix a nonempty spherical $T\subseteq S$.  For $\ell\in\widetilde T$
and $h\in A_T$ let
\[
 R_\ell(h)=U_\ell B_h.
\]
Expanding a word for $B_h$ gives one residual simple $U_\ell$ in
placements $(0,0)$ and shifted right projectives $Q_\beta$ with
$\beta\in\widetilde T$.  Minimize in the additive category generated
by these graded terms.  Its indecomposable endomorphism rings are local,
so the usual radical cancellation and uniqueness argument applies.
The simple cannot cancel with a projective, which has both top and
token socle.  Hence exactly one residual simple remains.  Denote by
$K^r_{\ell\beta}(h)$ the number of projective generators of type
$Q_\beta$ at total generator degree $r$, summed over both placements.

For representatives $z,z'$ of two candidates, set $h=z'z^{-1}$.
The identity
\[
 R_\ell(h)\otimes_ZE_H(z)\simeq U_\ell E_H(z')
\]
relates the transported right simple to the two left families.  Later
we will use their caps to bound the projective terms of $R_\ell(h)$.
The detection lemma below states exactly which vanishing will suffice:
once the two open sets agree and are denoted by $V$, vanishing of
$K^r_{\ell\beta}(h)$ for all $r$ and
$\ell,\beta\in\widetilde V$ forces $h\in A_{T\setminus V}$.
We establish this criterion, together with the invariance of the block
under changing representatives, before estimating it.
\begin{lemma}[Right duality and exterior invariance]
\label{two:duality}
For all $\ell,\beta\in\widetilde T$ and $r\in\mathbb Z$,
\begin{equation}
 K^r_{\ell\beta}(h)=K^{-2-r}_{\beta\ell}(h^{-1}).
 \label{two:duality-equation}
\end{equation}
For $J,J'\subseteq T$, left multiplication of $h$ by $A_{J'}$ leaves all rows outside
$\widetilde J'$ unchanged.  Right multiplication by $A_J$ leaves
all columns outside $\widetilde J$ unchanged.
\end{lemma}
\begin{proof}
First minimize the bimodule complex $B_h$.  Before minimization it
has one regular bimodule $Z$ at $(0,0)$, and all other terms are
shifts of $P_\alpha\otimes_{\kk}Q_\beta$, with indices in
$\widetilde T$; tensoring two such terms has the same form because
their middle tensor factor is a graded path space.

No indecomposable regular summand is isomorphic to a projective
bimodule of this form.  A connected-component regular bimodule is
indecomposable: a central idempotent would separate its vertices,
contrary to the nonzero arrows connecting them.  Modulo the left
and right radicals, its top contains the diagonal simple at every
vertex of that component.  The top of $P_\alpha\otimes Q_\beta$
contains only the pair $(\alpha,\beta)$.  The only remaining
possibility is a singleton component with $\alpha=\beta$; there the
regular bimodule has dimension two and the projective bimodule has
dimension four.  Shifts do not change this obstruction.  Thus all
regular summands survive, only at $(0,0)$.

A radical differential between projective bimodules is a sum of path
tensors with positive length on at least one side: a nonzero
scalar component with both lengths zero would be an invertible
entry and would have been canceled.  Tensoring by $U_\ell$ kills
positive left lengths; every surviving projective differential has
positive right length.  Maps involving the regular term specialize
to radical maps between the simple and a projective.  There is no
differential between residual simples, since they occur only in
cochain degree zero.  The specialized right complex is therefore
already minimal, and its counts are the bimodule counts with first
index $\ell$.

The Frobenius trace gives adjoint duality
\[
 \Hom_{Z\text{-left}}(M,Z)\cong\Hom_{\kk}(M,\kk)\langle2\rangle.
\]
It preserves additive radicals and sends
$P_\ell\otimes Q_\beta$ to
$(P_\beta\otimes Q_\ell)\langle-2\rangle$.  On complexes it reverses
cochain placements as well.  Thus $(c,u)$ becomes $(-c,-u-2)$.
The adjoint of the invertible complex $B_h$ is its inverse
$B_{h^{-1}}$, by Proposition~\ref{cat:action}.  Minimality and the
preceding specialization give \eqref{two:duality-equation}.

If $\ell\notin\widetilde J'$, every letter of $A_{J'}$ fixes
$U_\ell$: its inserted bimodule terms tensor to zero.  This proves
row invariance.  Apply \eqref{two:duality-equation} to
$(ha)^{-1}=a^{-1}h^{-1}$ to obtain column invariance for $a\in A_J$.
\end{proof}

To descend a cohomologically detected unfolded letter to a color,
we record the required finite-type greedy argument.  A simple positive
element is a reduced-word lift of an element of $W_T$.  By the
finite-type monoid facts in Theorem~\ref{cat:classical}, a
positive element $b$ has a greatest simple right divisor $t(b)$,
its right gcd with the Garside element.

\begin{lemma}[Greedy divisibility and unfolding]
\label{two:greedy}
For nonidentity simple factors $d_1,\ldots,d_k$, repeated extraction
of the greatest simple suffix gives $d_1\cdots d_k$ precisely when
\begin{equation}
 D_R(d_i)\subseteq D_L(d_{i+1})\qquad(1\le i<k).
 \label{two:greedy-condition}
\end{equation}
If an unfolded vertex right-divides the unfolding of a positive
element of $A_T$, its color right-divides the original element.
The analogous assertion holds for left divisors.
\end{lemma}
\begin{proof}
An atom $s$ right-divides a simple $d$ exactly when $s\in D_R(d)$,
and $sd$ is simple exactly when $s\notin D_L(d)$.  Failure of
\eqref{two:greedy-condition} therefore transfers a final atom of
$d_i$ to the beginning of $d_{i+1}$ and enlarges the simple suffix
at that extraction stage.  This proves necessity.

For sufficiency, induct on prefixes.  Suppose
$P=d_1\cdots d_{j-1}$ has greatest simple suffix $d_{j-1}$ and
write $t(Pd_j)=cd_j$, since $d_j$ is itself a simple suffix.  If
$c\ne1$, choose a final atom $s$ of $c$.  Then $sd_j$ is a factor
of a simple element, so it is simple.  Cancellation in
$Pd_j=bcd_j$ gives $P=bc$, whence $s$ right-divides $P$ and thus
$t(P)=d_{j-1}$.  This contradicts
$D_R(d_{j-1})\subseteq D_L(d_j)$.  Hence $c=1$, completing the
induction and the successive extraction proof.

Unfold a greedy factorization.  By Lemma~\ref{cat:unfolding},
its factors remain simple and their descent sets become full color
preimages.  Condition \eqref{two:greedy-condition} still holds, so
the unfolded factorization is greedy.  Any final unfolded atom
divides its last simple factor.  Descent reflection makes its color
a final divisor of the last original factor, hence of the original
element.  Word reversal preserves all positive braid relations and
interchanges left and right divisors, proving the left version.
\end{proof}

The following detector adapts the common-divisor cancellation used in
cohomological Garside detection; compare the proof of
\cite[Theorem~3.1]{BravThomas2011}.  Here the tests are right simples,
and the conclusion identifies a standard parabolic subgroup.

\begin{lemma}[Spherical parabolic detection]
\label{two:detection}
If $\varnothing\ne V\subseteq T$ and
\begin{equation}
 K^r_{\ell\beta}(h)=0
 \qquad(r\in\mathbb Z,\ \ell,\beta\in\widetilde V),
 \label{two:zero-block}
\end{equation}
then $h\in A_{T\setminus V}$.
\end{lemma}
\begin{proof}
Fix $\ell\in\widetilde V$.  A right projective has a nonzero map
to or from $U_\ell$ only when it is $Q_\ell$, via its top quotient
or token socle.  Equation \eqref{two:zero-block} excludes that
projective from $R_\ell(h)$, so its residual simple splits off as
a complex.  Tensoring with $B_h$ preserves indecomposability.
Therefore the other summand is zero in the homotopy category, and
$U_\ell B_h\simeq U_\ell$.

Write $h=h_1h_2^{-1}$ with $h_1,h_2$ positive, and set
$U_V=\bigoplus_{\ell\in\widetilde V}U_\ell$.  Then
\begin{equation}
 U_VB_{h_1}\simeq U_VB_{h_2}.
 \label{two:positive-equality}
\end{equation}
Use the maximum total test height on $\widetilde T$, which is zero
on $U_V$.  A positive word avoiding $V$ fixes $U_V$.  If it contains
a color of $V$, all operations before the first such color fix
$U_V$, and that first operation raises height by acting on a
highest test.  Positive monotonicity then keeps the height positive.
Thus final height zero is equivalent to avoiding $V$.

If the common height in \eqref{two:positive-equality} is positive,
choose a highest-test vertex.  The right version of
Lemma~\ref{lay:top-divisor}, valid for bounded tests including right
simples, extracts that vertex as a final atom of both unfolded
words.  Lemma~\ref{two:greedy} extracts a common final original
color $s$.  Write $h_i=g_i\sigma_s$ and tensor
\eqref{two:positive-equality} on the right by $B_{\sigma_s}^{-1}$.
The equality persists for $g_1,g_2$, while
$h_1h_2^{-1}=g_1g_2^{-1}$.  Their total positive length decreases
by two.  Iteration terminates at height zero, when both words avoid
$V$ and their fraction lies in $A_{T\setminus V}$.
\end{proof}

\subsection{Normalized right levels}

We next arrange the right-simple rows so that their highest layer can
be compared with a capped left layer. For projective terms, the tensor
calculation has three columns, indexed by the length $j=0,1,2$ of the
middle path in $Z$. A right generator of level $u$ and a left generator
of level $v$ contribute in total degree $u+v+j$.
Tensoring the residual simple, which stays in placements $(0,0)$,
with a left projective is nonzero only when their vertices agree,
and then has total degree $v$. We place these simple tensors in column one.
To retain the same degree formula, we assign the residual simple
\emph{normalized level} $-1$. This is a bookkeeping level; its actual
cochain and internal placements remain $(0,0)$.

Projective terms retain their usual generator level, the sum of their
cochain and internal placements.
A projective-to-simple differential entry is a top quotient with source
placements $(-1,0)$; a simple-to-projective entry is a token inclusion
with target placements $(1,-2)$. Their two endpoints consequently have
normalized level $-1$. Give these entries effective length one, as for
an ordinary arrow entry. Every entry then changes normalized level by
one minus its effective length. At each normalized level retain the
effective length-one entries. At a projective vertex, the resulting
right top is its coefficient space modulo incoming images, including
those from the residual simple.

If $r\geq-1$ bounds the normalized levels, then
\begin{equation}
 H^{r+2}(R_\ell(h)e_\beta)
 \text{ has dimension }\dim\Top_\beta(R_\ell(h)^r),
 \qquad
 H^t(R_\ell(h)e_\beta)=0\quad(t>r+2).
 \label{two:right-top}
\end{equation}
Indeed, only level-$r$ projective tokens can occur in degree $r+2$.
Their incoming images are the arrow-coefficient images, by the perfect
backtrack pairing, together with the simple's token inclusions when
$r=-1$. An additional incoming token differential would require a
level-$(r+1)$ source, which is absent. The residual simple itself has
degree zero, below $r+2$. The same argument applies to a common bound
for a finite family of rows.

\begin{lemma}[Negative right normalization]
\label{two:negative-bound}
A right inverse twist, followed by minimization, preserves an upper
normalized bound $r\ge-1$ for a complex with one residual simple
at $(0,0)$ and shifted right projectives.  At its operated vertex
$\alpha$ the tests satisfy
\begin{equation}
 (R D_\alpha)e_\alpha
       \simeq (Re_\alpha)\langle-2\rangle[-1].
 \label{two:negative-test}
\end{equation}
Both assertions hold for a full color twist as well.
\end{lemma}
\begin{proof}
For an individual unfolded vertex the inverse bimodule is
\[
 D_\alpha=[Z\longrightarrow
              (P_\alpha\otimes Q_\alpha)\langle-2\rangle],
\]
with its second term in cochain degree one.  A projective at
$(c,u)$ and a middle path of length $a$ insert $Q_\alpha$ at
$(c+1,u+a-2)$, of level $c+u+a-1$.  Only the insertions with
$a=2$ can exceed the old level, and they come from old
$Q_\alpha$ terms.

These offending terms can be canceled simultaneously, even in the
presence of the old differential.  In $\mathcal E=\operatorname{Tot}
(R D_\alpha)$ let $A$ be the sum of old $Q_\alpha$ terms in the
$Z$ column and $B$ their inserted length-two copies.  The latter
are at $(c+1,u)$, and the token--unit part of coevaluation gives
an invertible block
$\Phi=\pi_Bd|_A:A\to B$, consisting of signed scalar identities.
The old differential stays within a tensor column and does not
enter this block.  The subcomplex
\[
 \mathcal K=A\oplus d(A)
\]
is contractible: projection onto $B$ identifies $d(A)$ with $B$,
and its differential maps $A$ isomorphically to $d(A)$ and kills
the latter.  It is a direct summand as a bigraded module.
Writing $\mathcal E=A\oplus B\oplus W$, the quotient is homotopy
equivalent to $\mathcal E$ and has underlying terms $W$ with
\begin{equation}
 d'_W=d_{WW}-d_{WA}\Phi^{-1}d_{BW}.
 \label{two:pivot}
\end{equation}
Thus scalar entries induced among the inserted terms cause no
problem; their absence was not assumed.

The remaining projectives are the old terms of other vertices and
the length-zero and length-one insertions.  Their levels are,
respectively, the old level, the old level minus one, and the old
level.  The simple remains at normalized level $-1$; if its vertex
is $\alpha$, it also inserts $Q_\alpha$ at $(1,-2)$, again level
$-1$.  All remaining terms have level at most $r$.  Subsequent
minimization only removes summands and so preserves the bound.

The right-handed positive test identity follows from the same
token-cone cancellation as its left-handed version in
Lemma~\ref{lay:positive}.  Inverting it gives
\eqref{two:negative-test}, of total shift $-1$.  This is a bimodule
homotopy before tensoring, so it holds for simple inputs too.
Finally the vertices in one color are orthogonal, and their twists
commute; successive application proves the color statements.
\end{proof}

\subsection{Comparing two candidates}

We now compare two candidates with data
$(J,V,p,a,z)$ and $(J',V',p',a',z')$.  Initially choose $z$ capped
on $T$ by accessibility, while $z'$ need only have its prescribed
cap on $V'$.  Put $h=z'z^{-1}\in A_T$.  Use right rows
$R_\ell(h)$ for $\ell=(i,L)\in\widetilde V'$, with weights
$a'_iw_L$, and independently give the left member $E_{L'}(z)$
weight $w_{L'}$.  Since
\[
 R_\ell(h)\otimes_ZE_{L'}(z)\simeq U_\ell E_{L'}(z'),
\]
their weighted actual cohomology in total degree $d$ has dimension
\begin{equation}
 D_0\sum_{i\in V'}a'_i m_i^d(z').
 \label{two:actual-cohomology}
\end{equation}
Indeed tensoring a minimal left complex with $U_\ell$ kills its
radical differential and counts its generators at $\ell$.  Summing
over $L'$ first gives $w_Lm_i^d(z')$, and then the row sum gives
\eqref{two:actual-cohomology}.  It vanishes for $d>q$.

Let $r\ge-1$ be the highest normalized level of these right rows,
and let $N$ be their coefficient data at level $r$, retaining the
row index.  Put $M_{L'}=Y^q(E_{L'}(z))$.  The linear tensor complex
at the pair $(r,q)$ has columns of middle path lengths $0,1,2$,
with simple tensors assigned to column one.  Denote its cohomology
by $H_j(N,M)$ and its weighted dimensions by $\dim_{\wt}$.

\begin{lemma}[The highest tensor pair]
\label{two:tensor}
One has $H_2(N,M)=0$, and $H_1(N,M)$ survives in actual total
degree $r+q+1$.  At any projective vertex the token cokernel
surjects onto the tensor product of the right and left top quotients.
In particular a highest right top in a color $j\in J$ forces all
level-$q$ left-family tops in that color to vanish.
\end{lemma}
\begin{proof}
Filter the entire tensor by middle path length.  Radical projective
entries increase length by one or two.  A top quotient into the
simple takes a length-zero tensor to its assigned length one;
a token inclusion out of it takes assigned length one to length
two.  A left radical map tensored with a right simple vanishes.
Every differential therefore increases this index.

For a pair of normalized levels $(u,v)$ and middle length $j$,
actual total degree is $u+v+j$.  This also holds for simple tensors,
since their assigned right level and middle length are $-1$ and
$1$.  The length-one differential preserves the two levels.
After its cohomology the only remaining possible differential goes
from column zero to column two.  In particular every column-one
class survives.  This assertion can also be seen directly: writing
the differential as $d_1+d_2$, the cycle conditions are
$d_1u_0=0$ and $d_2u_0+d_1u_1=0$; the only secondary obstruction
is the class of $d_2u_0$ in column two.

Every nonzero tensor uses a left vertex in $T\cup\partial T$,
because right projectives are supported at vertices of $T$ and
only equal-vertex or adjacent paths exist.  Hence all relevant
left levels are at most $q$.  A secondary differential hitting
column two at $(r,q)$ would originate at a level pair of sum
$r+q+1$, since the differential raises total degree by one and
length by two.  Right levels are at most $r$ and relevant left
levels at most $q$, so no such source exists.  Therefore $H_2(N,M)$
also survives.  It lies in degree $r+q+2>q$, where
\eqref{two:actual-cohomology} is zero.

Project the token tensors at a vertex to the tensor of its two
coefficient spaces modulo incoming images.  Each length-one
boundary has an incoming coefficient on one side, so the projection
annihilates boundaries and is surjective onto the tensor of the top
quotients.  A nonzero right top thus forces every left top at that
vertex to vanish.  Family proportionality in
Proposition~\ref{cat:family} makes a left top occurring anywhere in
one color occur at every label in some member.  Thus the entire
left-family top in that color vanishes.
\end{proof}

\begin{lemma}[Removal of highest right tops]
\label{two:normalization}
Changing $z$ within $A_Jz$, while keeping its cap on $T$, one can
arrange that no highest right top occurs on $\widetilde J$.
\end{lemma}
\begin{proof}
For a highest right top in $j\in J$, Lemma~\ref{two:tensor}
licenses the cap-preserving positive move $z\mapsto\sigma_jz$.
It changes $h$ to $h\sigma_j^{-1}$ and applies the corresponding
inverse twist to the right rows.  Lemma~\ref{two:negative-bound}
shows that $r$ cannot increase.  Candidate data $p,a$ are unchanged,
and the tensor products still compute the same target at $z'$, so
the next move is licensed anew.

At a fixed $r$, the chosen negative right moves form a reduced word
in $W_J$.  Otherwise a right descent could be placed last by braid
moves.  All intermediate negative moves in that expression preserve
the upper bound $r$.  Its final acted-color test shifts down by one,
so its degree-$r+2$ test vanishes, contradicting the currently chosen
right top by \eqref{two:right-top}.  Sphericity bounds the length of
such a sequence.  Thus after finitely many moves either the relevant
tops disappear or $r$ decreases.  The residual simple bounds $r$
below by $-1$, so the procedure terminates.  If $J$ is empty no move
is required.
\end{proof}

\subsection{The weighted Euler estimate}

Use the representative in Lemma~\ref{two:normalization} and retain
$r,N,M$.  Put $m=m^q(z)$ and $e=m-p$.  The universal candidate
inequality verifies Proposition~\ref{cap:hull} for $A_Jz$: the cap
on $T\setminus J$ is fixed, so its joint $J$-cap is equivalent to
the additional cap on $T$.  Hence $e$ is a convex combination of
folded dimensions $s$ of subfamilies $\mathcal S\subseteq M$
supported on $\widetilde J$.

\begin{lemma}[Tensoring a subobject]
\label{two:subobject-tensor}
For each such subfamily,
\[
 H_0(N,\mathcal S)\lhook\joinrel\longrightarrow H_0(N,M),
 \qquad H_2(N,\mathcal S)=0.
\]
\end{lemma}
\begin{proof}
Column-zero cohomology is a kernel with no incoming boundaries,
so inclusion of its vector spaces gives the first assertion.

For the second work row by row and member by member, retaining
cochain placements.  At a projective vertex of $\widetilde J$,
each right coefficient is a sum of incoming images because its
top vanishes.  We use this fact inside the token cokernel.
Suppose an image at $\beta$ comes from a coefficient at $\gamma$
through a right arrow entry.  Choose a middle arrow dual to that
entry under the backtrack pairing and tensor it with the predecessor
coefficient and a left coefficient in $\mathcal S_\beta$.
Its differential has two possible token contributions: the chosen
right image paired with the unchanged left coefficient at $\beta$,
and the predecessor paired with the corresponding left-arrow image
at $\gamma$, with its cochain sign.  Non-backtrack products vanish;
dual arrow bases isolate the selected coefficient even for arrow
spaces of dimension greater than one.  Modulo boundaries, therefore,
the former token tensor equals the latter, up to a nonzero scalar
and sign.

The new left coefficient remains in $\mathcal S$ by the
subrepresentation property.  Each substitution decreases the right
cochain placement by one, so boundedness forces termination.  At a
projective vertex outside $\widetilde J$ the left coefficient is
zero.  An incoming image from the simple is a boundary directly:
the assigned length-one simple tensor maps to that token inclusion,
and its left differential is zero.  Thus every token tensor dies
in the cokernel, proving $H_2(N,\mathcal S)=0$.
\end{proof}

Define a nonnegative vector supported on $T$ by
\begin{equation}
 n_j=\frac1{D_0}\sum_H w_H
       \sum_{i\in V'}\sum_L a'_i w_L
                  K^r_{(i,L),(j,H)}(h).
 \label{two:right-weight}
\end{equation}
For a left family of weighted dimension $\phi y$, its weighted
length Euler characteristic with $N$, divided by $D_0$, is
\begin{equation}
 \chi(n,y)=n^{\mathsf T}Cy
           -\mathbf1_{\{r=-1\}}(a')^{\mathsf T}y_T.
 \label{two:euler-formula}
\end{equation}
For completeness, let
$N_\beta^{\wt}=\sum_{i\in V',L}a'_iw_LK^r_{(i,L),\beta}(h)$.
The projective terms contribute
\[
 \sum_{\beta,\gamma}N_\beta^{\wt}\widetilde C_{\beta\gamma}
              (\phi y)_\gamma
 =\sum_{j,H}N_{(j,H)}^{\wt}w_H(Cy)_j
 =D_0n^{\mathsf T}Cy.
\]
The unit and token contribute plus signs, and the middle arrows
minus signs.  We used $\widetilde C\phi=\phi C$.  At $r=-1$ the
residual simples contribute in column one
$-\sum_{i,L}a'_iw_L(w_Ly_i)=-D_0(a')^{\mathsf T}y_T$.
At higher $r$ there is no residual simple.  This proves
\eqref{two:euler-formula}, including its normalization.

\begin{lemma}[The two-point Euler estimate]
\label{two:euler}
At the normalized representative,
\begin{equation}
 \frac{\dim_{\wt}H_1(N,M)}{D_0}
 \ge n_T^{\mathsf T}a+
                \mathbf1_{\{r=-1\}}(a')^{\mathsf T}p_T.
 \label{two:estimate}
\end{equation}
\end{lemma}
\begin{proof}
For each hull vector $s$, Lemma~\ref{two:subobject-tensor} gives
\[
 \chi(n,s)=\frac{\dim_{\wt}H_0(N,\mathcal S)
                     -\dim_{\wt}H_1(N,\mathcal S)}{D_0}
       \le\frac{\dim_{\wt}H_0(N,M)}{D_0}.
\]
Linearity gives the same upper bound for $\chi(n,e)$, since $e$
is a convex combination of these vectors.  Lemma~\ref{two:tensor}
gives
$\chi(n,m)=(\dim_{\wt}H_0(N,M)-\dim_{\wt}H_1(N,M))/D_0$.
Subtracting yields
\[
 \frac{\dim_{\wt}H_1(N,M)}{D_0}
       \ge\chi(n,e)-\chi(n,m)=-\chi(n,p).
\]
Now $n$ is supported on $T$ and $(Cp)_T=-a$, so
\eqref{two:euler-formula} is exactly \eqref{two:estimate}.
\end{proof}

\begin{proof}[Proof of Proposition~\ref{two:uniqueness}]
Normalize the first representative as above.  If $r>-1$, its highest
right layer has nonzero projective coefficients and no residual
simple.  Some coefficient lies in $\widetilde V$: otherwise all
would lie in $\widetilde J$, where their tops vanish, and iterating
incoming images toward smaller cochain placements would force the
entire layer to vanish.  All row weights are positive, so
$n_T^{\mathsf T}a>0$.  Equation \eqref{two:estimate} gives nonzero
middle cohomology surviving in degree $r+q+1>q$, contradicting
\eqref{two:actual-cohomology}.  Therefore $r=-1$.

The surviving middle part now lies in degree $q$ and has dimension
at most the actual degree-$q$ cohomology.  Since $a'$ is supported
on $V'$, where $p'_i=m_i^q(z')$, we obtain
\begin{equation}
 (a')^{\mathsf T}(p'_T-p_T)\ge n_T^{\mathsf T}a\ge0.
 \label{two:forward}
\end{equation}
Interchanging the candidates gives
\begin{equation}
 a^{\mathsf T}(p_T-p'_T)\ge0.
 \label{two:reverse}
\end{equation}
The two normalizations may use different representatives, but their
harmonic vectors and open-color vectors belong to the candidate
cosets and are unchanged.  The exterior values of $p,p'$ outside
$T$ agree.  With $\delta=p'_T-p_T$ this gives $a'-a=-C_T\delta$.
Adding \eqref{two:forward} and \eqref{two:reverse} yields
$-\delta^{\mathsf T}C_T\delta\ge0$.  Positive definiteness forces
$p=p'$, hence $a=a'$, $V=V'$ and $J=J'$, with equality throughout
and $n_T^{\mathsf T}a=0$.

It remains to identify the cosets.  Fix any original comparison
$h_0=z'z^{-1}$.  Replacing the two representatives by $bz$ and
$cz'$, with $b\in A_J$, $c\in A_{J'}$, replaces $h_0$ by
$ch_0b^{-1}$.  Lemma~\ref{two:duality} preserves its block with
rows in $\widetilde V'$ and columns in $\widetilde V$.  These are
now the same set $\widetilde V$.

The forward comparison has no right generators above level $-1$.
At level $-1$, equality $n_T^{\mathsf T}a=0$ and
\eqref{two:right-weight} imply
$K^{-1}_{\ell\beta}=0$ for $\ell,\beta\in\widetilde V$, since
every coefficient weighting such a count is strictly positive.
Block invariance therefore gives the same zero block for $h_0$
in all degrees $r\ge-1$.  The reverse comparison gives its zero
block for $h_0^{-1}$ in all degrees $r\ge-1$, regardless of its
chosen representatives.  Duality \eqref{two:duality-equation}
transposes this into vanishing for $h_0$ in all degrees $r\le-1$.
Together the two ranges cover every integer.  By
Lemma~\ref{two:detection}, $h_0\in A_{T\setminus V}=A_J$.
Thus $A_Jz'=A_Jz$, proving uniqueness of the typed candidate.
\end{proof}
